\documentclass[multi=page, border=15pt, varwidth=16cm]{standalone}
\usepackage{ctex}      % 提供中文支持
\usepackage{tabularx}  % 自动调整列宽
\usepackage{booktabs}  % 优雅的三线表支持
\usepackage{xcolor}
\usepackage{colortbl}

% 统一的深灰蓝表头色
\definecolor{headerbg}{RGB}{52, 73, 94}
\definecolor{headertext}{RGB}{255, 255, 255}

% 自定义表头文本样式命令
\newcommand{\thnode}[1]{\textcolor{headertext}{\bfseries #1}}

\begin{document}
\sffamily
\renewcommand{\arraystretch}{1.5}

% ================= 图片：parted 核心指令集 =================
\begin{page}
	\centering
	% \textbf{\Large parted 磁盘管理核心指令集} \\ \vspace{0.8em}

	\begin{tabularx}{\linewidth}{>{\bfseries\color{blue!70!black}}l X >{\bfseries\color{blue!70!black}}l X >{\bfseries\color{blue!70!black}}l X}
		\toprule
		\rowcolor{headerbg}
		\thnode{指令} & \thnode{说明}      & \thnode{指令} & \thnode{说明}    & \thnode{指令} & \thnode{说明}   \\
		\midrule
		% 第一行：基础信息
		print       & \textbf{显示分区表信息} & help        & 显示帮助信息         & select      & 切换目标设备        \\
		% 第二行：初始化与创建
		mklabel     & \textbf{创建分区表标签} & mkpart      & \textbf{创建新分区} & rm          & 删除指定分区        \\
		% 第三行：调整与单位
		resizepart  & 调整分区大小           & unit        & 更改显示单位         & align-check & 检查分区对齐        \\
		% 第四行：退出
		name        & 设置分区名称           & toggle      & 切换分区标志         & quit        & \textbf{退出程序} \\
		\bottomrule
		\multicolumn{6}{p{\linewidth}}{\small \color{gray}\textit{* 注：\texttt{mklabel} 可选 \texttt{gpt} 或 \texttt{msdos}(MBR)；\texttt{mkpart} 建议起始位置设为 \texttt{1MiB} 以确保 4K 对齐}}
	\end{tabularx}

	% \vspace{1.5em}

	% 实战命令示例
	\begin{tabularx}{\linewidth}{>{\columncolor{headerbg!10}}l X}
		\toprule
		\textbf{实战流} & \textbf{指令示例}                                                \\
		\midrule
		\textbf{快速分区}  & \texttt{parted /dev/sdb mklabel gpt mkpart primary 1MiB 100\%} \\
		\textbf{无损扩容}  & \texttt{parted /dev/sdb resizepart 1 5GiB} (将 1 号分区扩大至 5GiB)   \\
		\textbf{对齐检查}  & \texttt{parted /dev/sdb align-check optimal 1} (检查 1 号分区是否对齐)  \\
		\bottomrule
	\end{tabularx}
\end{page}

\end{document}